\section{Validation}
\label{sec:validation}

\subsection{Validation strategy}
\label{sec:validation-strategy}

No World Cup data exists before the tournament starts, so we validate
on the closest available proxy: the English Premier League fantasy
game, whose scoring rules are nearly identical to those of the FIFA
World Cup Fantasy game (the two games share a common heritage). A
model that ranks players well on EPL data should rank players well on
tournament data, subject to the distribution-shift caveats analysed
in Section~\ref{sec:analysis}.

We hold out gameweeks 30--38 of the 2024--25 season (the last nine
gameweeks) and train on all earlier rows of the three-season corpus
of Section~\ref{sec:data-fpl} (2022--23, 2023--24, and 2024--25
through gameweek 29). Held-out error is the per-position root mean
squared error (RMSE) between the predicted score and the realised
\texttt{total\_points}. All three backends of
Section~\ref{sec:methodology} score exactly the same held-out rows,
so the comparison is paired; the per-position row counts appear in
Table~\ref{tab:holdout-rmse}.

\subsection{Held-out results and backend routing}
\label{sec:validation-holdout}

Table~\ref{tab:holdout-rmse} and Figure~\ref{fig:holdout-rmse} report
the held-out RMSE. No single backend wins everywhere; the winners
split cleanly by position.

\begin{itemize}
\item \textbf{Goalkeepers: Poisson} (\RmseGkPoisson{} against
  \RmseGkGbm{} for the GBM). Goalkeeper scoring is dominated by
  appearance and clean-sheet points, and the clean-sheet probability
  is exact under the Poisson goal model, so the structural backend
  captures most of the available signal directly.
\item \textbf{Defenders: heuristic} (\RmseDefHeuristic{}). The
  price-coefficient-times-matchup formula tracks defender outcomes
  better than the GBM at this data scale.
\item \textbf{Midfielders and forwards: GBM} (\RmseMidGbm{} and
  \RmseFwdGbm{}). Individual variance is largest at the attacking
  positions, and the GBM has room to learn non-linear interactions
  that a compact hand-written formula cannot capture.
\end{itemize}

This routing is not merely descriptive: it defines the per-position
ensemble backend, which runs the component models and keeps, for
each position, the prediction of the routed winner. A sixth candidate
GBM feature, the team Elo difference, was evaluated under the same
protocol and rejected; we report that negative result in
Section~\ref{sec:negative-results}.

\subsection{Two GBM iterations}
\label{sec:validation-gbm}

The GBM column of Table~\ref{tab:holdout-rmse} is the second
iteration of the model. The first (v1) trained on a single EPL season
with LightGBM defaults, 31 leaves and 400 boosting rounds
\citep{ke2017lightgbm}; the second (v2) trains on all three seasons
with a lighter sweep-tuned configuration, 15 leaves and 200 rounds,
and improves the held-out RMSE at every position. The defect that
motivated the sweep was visible in squad space before it was visible
in aggregate error: v1 squads were aggressive and implausible, with
clusters of cheap defenders from unfancied squads promoted by a
handful of training rows, a classic signature of too many leaves on
too little data. The v2 squads are conventional. We treat this as a
practical reminder that inspecting the downstream decisions catches
failure modes that aggregate metrics dilute.

\subsection{Determinism and the noise floor}
\label{sec:validation-determinism}

Feature evaluation at these effect sizes requires reproducible
training. Before we pinned the random state, identical code on
identical data produced per-position RMSE values that varied by
roughly $\pm 0.05$ between runs, driven by LightGBM's bagging
randomness. That wobble was large enough to mask the Elo-feature
comparison above: the first attempt produced mixed deltas
indistinguishable from noise. Setting \texttt{seed},
\texttt{bagging\_seed}, \texttt{feature\_fraction\_seed}, and the
deterministic flag collapsed the run-to-run variance to zero and made
the comparison honest. We now treat the pinned seed as part of the
training contract: every RMSE reported in this section is exactly
reproducible.

\subsection{Leak-free walk-forward validation on tournament data}
\label{sec:validation-walkforward}

The EPL hold-out cannot detect EPL-to-tournament distribution shift,
and it cannot evaluate features that only exist once tournament
rounds complete. Mid-tournament we therefore added a second protocol:
leak-free walk-forward validation on realised World Cup rounds. For
each completed round $k \ge 2$, the GBM trains on all EPL rows plus
the WC rounds strictly before $k$, then predicts round $k$; round 1
is excluded because no prior tournament round exists from which to
compute form. Squared errors are pooled across held-out rounds 2--6
to give the per-position and overall RMSE of
Table~\ref{tab:walkforward}; Figure~\ref{fig:walkforward} shows the
per-round traces. Six configurations are compared; the last two draw
on a community-maintained dataset of real match statistics
\citep{mominul2026dataset} that became part of the pipeline before
the final round.

\begin{itemize}
\item \textbf{A: EPL only, no form.} The shipped v2 model: five price
  and team-strength features, trained on EPL rows only. This
  configuration is form-blind; no feature tells it how a player has
  actually performed in the tournament.
\item \textbf{B: EPL only, plus form.} Adds \texttt{form\_lag}, the
  trailing mean of the player's realised fantasy points over the
  previous three matches, shifted by one match so a row's label never
  enters its own feature. One definition is shared by EPL training,
  WC label extraction, and WC inference; a player's first match
  yields a missing value, which LightGBM handles natively.
\item \textbf{C: EPL plus WC, plus form.} As B, but the training set
  also includes every realised WC round before the held-out round,
  recalibrating the output scale to the tournament.
\item \textbf{D: EPL plus WC, proxy minutes.} As C, plus a lagged
  start-rate feature (participation proxied by positive scores) and a
  team goals-conceded form feature.
\item \textbf{E: EPL plus WC, real xG form.} As C, plus the trailing
  mean of the team's real expected goals for and against over its
  previous three matches, computed leak-free from the community match
  dataset. The columns exist only for tournament rows; EPL training
  rows carry missing values, which LightGBM handles natively.
\item \textbf{F: EPL plus WC, real minutes.} As C, plus real lagged
  start rates and minutes shares from the same dataset's lineup
  sheets, testing whether the failure of configuration D was the
  proxy or the signal.
\end{itemize}

The pooled ordering is unambiguous. RMSE falls from
\WalkforwardPooledA{} (A) to \WalkforwardPooledB{} (B) to
\WalkforwardPooledC{} (C), improving at every position at each step,
with the largest single gain at goalkeeper once tournament labels
enter the training set. Configuration D edges C at goalkeeper and
defender but is worse at midfielder and forward and pools slightly
worse (\WalkforwardPooledD{}). Configuration E improves on C at
\emph{every} position and wins the pool at \WalkforwardPooledE{},
the only feature candidate of the tournament to do so; it was
deployed as the fourth-generation model two days before the final
under a pre-registered rule (deploy only on a pooled improvement of
at least $0.01$ with no more than one position regressing).
Configuration F cleared the pooled threshold only marginally
(\WalkforwardPooledF{}), regressed forwards, and lost to E in the
same run, so real minutes stayed out of the model; we return to it
in Section~\ref{sec:negative-results}. On the round-of-16 feature
set the form-retrained model ranks the elite in-form forwards,
including Mbapp\'e, Kane, and Haaland, at the top of its
predictions, where the form-blind model had buried them below
midfielders from strong teams.

The EPL hold-out of Table~\ref{tab:holdout-rmse} is kept at the
version that defined the ensemble routing. Re-run under the deployed
fourth-generation feature set, the GBM improves at every position on
the same held-out rows (goalkeeper \RmseGkGbmVFour{}, defender
\RmseDefGbmVFour{}, midfielder \RmseMidGbmVFour{}, forward
\RmseFwdGbmVFour{}) while leaving the routing unchanged: the Poisson
backend still wins goalkeepers and the heuristic still edges
defenders.

Because these are genuinely out-of-sample tournament rounds, we take
Table~\ref{tab:walkforward} as the leak-free evidence for the
form-retrained model. The squad-level backend comparison of
Section~\ref{sec:live-results} is partly in-sample for the retrained
GBM and the ensemble, since the deployed model has seen the labels of
the rounds scored there; those columns are upper bounds, not
out-of-sample estimates.

\subsection{What the validation does not measure}
\label{sec:validation-limits}

Held-out EPL RMSE is a useful proxy for prediction quality, but
several phenomena specific to international tournament play are
invisible to it.

\begin{itemize}
\item \textbf{Star international players are absent from training.}
  Messi, Mbapp\'e, Modri\'c, Neymar, and others never appear in the
  EPL feature set, so the EPL-trained GBM scores them on price alone
  and undervalues them systematically. The walk-forward retrain above
  mitigates this once realised tournament rounds accumulate.
\item \textbf{Manager rotation behaviour differs.} EPL coaches rotate
  on fixture congestion; international coaches rotate on group-stage
  clinching. The model has no signal for either.
\item \textbf{Knockout-stage variance is higher.} Single-leg
  elimination games carry different team incentives than round-robin
  group games, and every EPL training row comes from a 38-game
  round-robin schedule.
\end{itemize}

These effects surface in the live-tournament results
(Section~\ref{sec:live-results}) and in the analysis
(Section~\ref{sec:analysis}), not in the held-out RMSE.

\begin{table}[t]
  \centering
  % GENERATED by scripts/paper_tables.py -- do not edit
\begin{tabular}{lrrrr}
\toprule
Position & $n$ & Heuristic & Poisson & GBM \\
\midrule
GK & 180 & 2.698 & \textbf{2.491} & 2.650 \\
DEF & 910 & \textbf{3.141} & 3.400 & 3.187 \\
MID & 1228 & 2.874 & 4.371 & \textbf{2.767} \\
FWD & 368 & 3.515 & 4.560 & \textbf{3.153} \\
\bottomrule
\end{tabular}

  \caption{Held-out EPL RMSE per position and backend. Bold marks the
    per-position winner; this routing defines the ensemble.
    \figsource{fifa\_fantasy.training.validate\_main $\rightarrow$
    data/training/validation\_report.json}}
  \label{tab:holdout-rmse}
\end{table}

\begin{figure}[t]
  \centering
  \includegraphics[width=0.85\textwidth]{fig_holdout_rmse.pdf}
  \caption{Held-out RMSE by position and backend.
    \figsource{fifa\_fantasy.report.figures $\rightarrow$
    data/training/validation\_report.json}}
  \label{fig:holdout-rmse}
\end{figure}

\begin{table}[t]
  \centering
  % GENERATED by scripts/paper_tables.py -- do not edit
\begin{tabular}{lrrrrr}
\toprule
Config & GK & DEF & MID & FWD & All \\
\midrule
A: EPL, no form & 3.517 & 3.122 & 3.063 & 3.185 & 3.164 \\
B: EPL + form & 3.350 & 2.962 & 2.776 & 2.984 & 2.952 \\
C: EPL+WC + form & 2.450 & 2.889 & \textbf{2.640} & \textbf{2.757} & \textbf{2.724} \\
D: EPL+WC, all features & \textbf{2.445} & \textbf{2.877} & 2.659 & 2.800 & 2.735 \\
\bottomrule
\end{tabular}

  \caption{Leak-free walk-forward RMSE on WC rounds 2--6, pooled per
    feature configuration. Bold marks the per-column best.
    \figsource{scripts/wc\_forward\_validation.py $\rightarrow$
    data/evaluation/wc\_forward\_validation.json}}
  \label{tab:walkforward}
\end{table}

\begin{figure}[t]
  \centering
  \includegraphics[width=0.85\textwidth]{fig_walkforward.pdf}
  \caption{Walk-forward RMSE per round and configuration.
    \figsource{scripts/wc\_forward\_validation.py}}
  \label{fig:walkforward}
\end{figure}
